\chapter{答案、提示与进一步阅读}

翻到这一页，你可能想核对某道练一练的答案。先说清楚这个附录\textbf{不是什么}：它不是一本标准答案册。全书八十七章的练一练，大多不是“默写某个名词”式的题目，而是让你画一张粒子图、能量图、物质流图或信息流图——这类题目的目的，是检查你有没有真的把模型装进脑子里。画图题没有唯一正确的画法，只有\textbf{违反规则的错法}。所以这个附录先做一件更要紧的事：教你一套“自己当判卷老师”的方法。学会之后，绝大多数练习你都能自己判断对错，不必等别人公布答案。

然后，我们挑六个风格差异最大的样章——第 10、17、28、49、69、73 章——给出每道题的完整提示与解答思路，示范这套自检方法怎么用。查符号和配平方法请去附录 A，周期表趋势请去附录 B，分子模型画法请去附录 C，生物学术语请去附录 D，这里都不重复。最后，按十一篇分别列出进一步阅读的建议。

\section{怎样自己检查答案：四张图的自检清单}

每画完一张图，对照相应清单逐条打勾。任何一条过不了，答案就是错的，不管画得多漂亮。

\subsection{粒子图自检清单}

粒子图把镜头推到原子、分子、离子层面，画出“谁在哪儿、在干什么”。检查五条：

\begin{itemize}[itemsep=2pt]
\item \textbf{种类对吗？}你画的小球到底是原子、分子还是离子？有没有把“元素”当成会跑的小球？（元素是种类，不是个体；会反应的是具体的原子和分子。）
\item \textbf{数量守恒吗？}把图左和图右每种原子的个数数一数，必须相等。化学反应不消灭原子，只重新组合——你的图里有没有原子凭空出现或消失？
\item \textbf{电荷守恒吗？}离子要标出电荷，且两边总电荷相等。钠丢给氯的那个电子，在你的图里必须找得到去向。
\item \textbf{结构合理吗？}分子的连接方式符合各原子的成键能力吗？有没有让氢长出两只胳膊？
\item \textbf{注明了吗？}图旁有没有写“这是人为的表示法，颜色和大小不代表真实外观”？谁把示意图当成照片，谁就会被示意图骗。
\end{itemize}

\subsection{能量图自检清单}

能量图回答“变化为什么能发生、要跨过多高的坎”。检查五条：

\begin{itemize}[itemsep=2pt]
\item \textbf{两根轴画对了吗？}纵轴是能量，横轴是“反应进程”或“原子间距”——不是时间。能量图不告诉你快慢。
\item \textbf{坡的方向对吗？}断键必须\textbf{先爬坡}（吸收能量），成键才\textbf{下坡}（放出能量）。凡是画成“断开化学键释放能量”的图，一律判错——这是全书反复划的红线。
\item \textbf{两根箭头分开了吗？}活化能（谷底到山顶的高度）和 $\Delta G$（起点与终点的高差）是两根不同的箭头。改变化学速率动的是前者，决定方向和平衡位置的是后者。
\item \textbf{催化剂画对了吗？}加催化剂或酶，只把山顶削矮（换一条路），起点和终点一动不动。谁把 $\Delta G$ 箭头也改短了，谁就中了“酶是能量来源”的招。
\item \textbf{方向和快慢分清了吗？}$\Delta G<0$ 只说“能发生”，不说“会快发生”。速度要看山有多高、温度多高，那是另一张图的事。
\end{itemize}

\subsection{物质流图自检清单}

物质流图追踪物质和能量在体系之间怎么进出、怎么循环。检查四条：

\begin{itemize}[itemsep=2pt]
\item \textbf{每个方框平账吗？}对每个“库”问一句：进的等于出的加积累的吗？有没有物质从方框里凭空长出来？
\item \textbf{每个箭头有推动力吗？}随便指一个箭头，你能说出它靠什么驱动吗——扩散、泵、化学反应、食物链？说不出推动力的箭头是装饰品。
\item \textbf{物质流和能量流分开了吗？}这是最容易错的一条：物质可以循环，能量只能\textbf{单向流过}——从阳光到食物链再到热，散掉就回不来。画成“能量循环”的图直接判错。
\item \textbf{循环闭合了吗？}叫“循环”的图，箭头必须能绕回起点。缺一段就不是循环，是有去无回的线。
\end{itemize}

\subsection{信息流图自检清单}

信息流图追踪遗传信息怎么被抄写、换码、执行，是第八到第十一篇的主力图。检查五条：

\begin{itemize}[itemsep=2pt]
\item \textbf{大方向对吗？}主线是 DNA $\rightarrow$ RNA $\rightarrow$ 蛋白质。复制和转录不换语言，都是核苷酸字母；只有翻译才真正“换码”。把换码标在转录上的，回去重读第 68、69 章。
\item \textbf{执行者标了吗？}复制靠 DNA 聚合酶，转录靠 RNA 聚合酶，翻译靠核糖体加 tRNA。信息不会自己流动，每一步都有分子机器。
\item \textbf{能量花在哪标了吗？}复制、装货、进位都要付能量账。不标能量去向，就等于假装信息处理免费。
\item \textbf{中间环节省略了吗？}从基因到性状之间至少隔着蛋白质、细胞、个体、环境好几层。一个箭头直接从基因指向性状的图，是把第 73 章整章跳过了。
\item \textbf{“序列”和“使用方式”分清了吗？}同一段 DNA，读不读、何时读、读多少，是另一层信息（第 70、74 章）。图里给调控留位置了吗？
\end{itemize}

把四张清单用熟，你就拥有了一件比答案更值钱的东西：\textbf{判断力}。下面六个样章的解答，就是这套清单的实战示范。每道题先给一句提示（合上书，看提示能不能让你自己走出来），再给完整的解答思路。

\section{样章一：第 10 章《周期表不是元素通讯录》}

本章的核心模型只有一句话：\textbf{最外层电子数决定化学性格}。所有题目都是它的应用。

\begin{enumerate}[itemsep=3pt]
\item （镁与氧相遇，谁失电子、谁得电子？）\textbf{提示：}别背“金属失电子”，算两笔账——丢几个划算，还是抢几个划算。\textbf{思路：}镁最外层 2 个电子：丢 2 个就露出住满的第二层，抢 6 个才能填满第三层——丢 2 个省力得多，变成 $\chem{Mg^{2+}}$；氧最外层 6 个：抢 2 个就住满，丢 6 个代价太大，变成 $\chem{O^{2-}}$。生成物 $\chem{MgO}$。判据不是“原子想要 8 个电子”，而是\textbf{哪条路让体系总能量降得多}；自检：你的解释里有没有偷偷给原子安上愿望？
\item （氢可燃、氦安全，用楼层规则解释。）\textbf{提示：}比较两种原子“重新排布电子”有没有赚头。\textbf{思路：}氢最外层只有 1 个电子，和氧共享成键后体系能量大幅下降，所以燃烧剧烈；氦的第一层只有 2 个名额且已住满，任何重组都不能再降低能量，于是对一切反应无动于衷。可不可燃是电子排布问题，和“氢比空气轻”无关——自检：你有没有把物理性质和化学性格混为一谈？
\item （门捷列夫凭什么把碲排在碘前面？）\textbf{提示：}他相信“性格”胜过相信“数字”。\textbf{思路：}依据是化学性质的周期性——碘明显属于卤素一家，碲属于氧硫一家，他宁可怀疑原子量测错。后来证明他对的关键发现是：真正决定顺序的是原子核里的质子数（原子序数），由莫塞莱用 X 射线实验测定；碲的质子数就是比碘少一个，原子量反常是同位素比例造成的。
\item （画钠和氯的电子入住图。）\textbf{提示：}先数电子总数，再从一楼按容量安排。\textbf{思路：}钠 11 个电子，住法 2-8-1——三楼刚开张，只 1 个客人；氯 17 个电子，住法 2-8-7——三楼差 1 个满员。用箭头标出：钠的三楼“急于退租”，氯的三楼“差一间满员”。图旁必须注明这是表示法，电子并不真的住在格子里——粒子图清单第五条。
\item （同族越往上，抢电子能力怎么变？）\textbf{提示：}把第 1 族“楼层越高越容易丢电子”的逻辑镜像过来。\textbf{思路：}丢电子看外层电子被吸得多松，抢电子看核对外来电子吸得多紧。同族越往上，最外层离核越近、屏蔽越少，核的吸引力越“够得着”外来电子，所以抢电子越强：氟强于氯，氯强于溴。自检：你的解释是否落在“吸引力与距离、屏蔽的拉锯”上，而不是背箭头方向？
\end{enumerate}

\section{样章二：第 17 章《化学键不是原子之间的小棍》}

本章的核心是一张图：两个原子靠近时，体系能量随距离先降后升，谷底就是化学键。

\begin{enumerate}[itemsep=3pt]
\item （画能量—距离曲线并解释“自发停在谷底”。）\textbf{提示：}先定零点，再想“近了为什么会涨回去”。\textbf{思路：}右端（相距无穷远）能量取零；靠近时核对电子的吸引占上风，曲线下降；谷底标出键长和键能；再近，核与核、电子云与电子云的排斥暴涨，曲线陡然翘起。“停在谷底”不是原子喜欢那儿，而是\textbf{任何偏离都会升高体系能量}——像球滚到碗底。条件：成键释放的多余能量必须能被带走（碰撞传出、发光放热），否则两个原子还会弹开。
\item （$\chem{H_2} + \chem{Cl_2} \rightarrow \chem{2HCl}$ 放热还是吸热？）\textbf{提示：}断键付钱，成键收钱，轧账。\textbf{思路：}断 1 个 $\chem{H-H}$（约 \ud{436}{kJ\cdot mol^{-1}}）加 1 个 $\chem{Cl-Cl}$（约 \ud{243}{kJ\cdot mol^{-1}}），共支付约 \ud{679}{kJ\cdot mol^{-1}}；成 2 个 $\chem{H-Cl}$（各 \ud{431}{kJ\cdot mol^{-1}}），共收回 \ud{862}{kJ\cdot mol^{-1}}。收回大于支付，净放约 \ud{183}{kJ\cdot mol^{-1}}，放热。若章内键能表数值略有出入，以表格为准，\textbf{方法不变}：支付端永远在断键一侧。
\item （“炸药存了很多高能的键，爆炸时断开键放出能量”错在哪？）\textbf{提示：}数出两处方向性错误，都用谷底语言重说。\textbf{思路：}错误一，\textbf{能量不存在键里}——断键永远吸能，不存在“断开键放出能量”这种事；错误二，\textbf{方向恰好相反}——爆炸放能是因为生成物（如 $\chem{N_2}$、$\chem{CO_2}$、$\chem{H_2O}$）的键更强、谷底更深，成键收回远超断键支付。正确说法：炸药分子是住在浅谷里的原子，一旦翻过小坡重排成深谷里的组合，落差就以热的形式冲出来。
\item （氮气多而懒、氧气少而凶，用键能解释。）\textbf{提示：}查三键和双键的键能差距。\textbf{思路：}$\chem{N \equiv N}$ 三键键能约 \ud{946}{kJ\cdot mol^{-1}}，拆开要付天价，活化能高得让大多数反应没机会启动；$\chem{O=O}$ 双键约 \ud{498}{kJ\cdot mol^{-1}}，便宜得多。所以占空气五分之四的氮气像个惰性观众。对你的意义：正因为氮气把氧气“稀释”了、自己又不掺和，纸张木头才不会见火星就失控。
\item （$\chem{SF_6}$ 和 $\chem{BF_3}$ 为什么不满足八隅体也稳定？）\textbf{提示：}口诀和法律，谁大？\textbf{思路：}八隅体是经验口诀，能量才是法律。硫在第三周期，“楼层”更大，周围容纳 12 个电子恰是能量最低的排法；硼在 $\chem{BF_3}$ 里 6 个电子的排法已足够低，再去“凑八”反而要付能量。口诀是能量规律的速记符号，好用但处处有例外（本章还给了 $\chem{NO}$、$\chem{XeF_4}$ 两组反例）；判断分子稳不稳，永远问“体系能量降了没有”。
\end{enumerate}

\section{样章三：第 28 章《能量下降还不够：熵与自由能》}

本章的核心是两条：方向由 $\Delta G = \Delta H - T\Delta S$ 裁决，熵的本质是“排法的多少”。

\begin{enumerate}[itemsep=3pt]
\item （画冰袋溶解的粒子图，标热量流向和熵增熵减。）\textbf{提示：}体系和环境分开标注。\textbf{思路：}左边画整齐的硝酸铵晶格加独立的水分子；右边画分散开的、被水分子包围的 $\chem{NH_4^+}$ 和 $\chem{NO_3^-}$。热量箭头从环境（你的手）指向体系，所以手觉得凉。熵的标注：体系从有序晶格变成分散离子，\textbf{熵增}；环境失热，\textbf{熵减}；但体系熵增更大，总账为正，$\Delta G<0$——吸热也能自发，这正是本章的招牌反例。
\item （墨水为什么自己散开、永不聚回？倒放视频违反力学吗？）\textbf{提示：}数一数两种状态的“排法”各有多少。\textbf{思路：}均匀分散的排法比聚成一滴多出天文数字倍——分子每次随机走动，走向均匀的概率压倒性地大。散开不是被力量“拉”开的，是被概率淹没的。倒放视频里，每次碰撞都不违反牛顿定律——力学对时间反演一视同仁；它违反的是统计：那样的巧合在宇宙年龄里也等不到一次。第二定律不是“禁止”，是“概率小到等于禁止”。
\item （为什么冷冻室结冰、冷藏室不结冰？）\textbf{提示：}$\Delta H$ 和 $T\Delta S$ 谁主导，由 $T$ 决定。\textbf{思路：}结冰 $\Delta H<0$（有利）但 $\Delta S<0$（不利）。温度低时 $T\Delta S$ 项很小，负的 $\Delta H$ 说了算，$\Delta G<0$，结冰——$-18\,^{\circ}\mathrm{C}$ 的冷冻室正是如此；温度高时熵减的代价盖过放热的好处，$\Delta G>0$，绝不结冰——$10\,^{\circ}\mathrm{C}$ 的冷藏室如此。“扯平”的条件是 $\Delta G = 0$，即 $T = \Delta H / \Delta S$，常压下就是 $0\,^{\circ}\mathrm{C}$。
\item （把四个过程按体系熵增熵减分类。）\textbf{提示：}盯住“粒子自由度”和“气体分子数”。\textbf{思路：}香水挥发，液变气，排法暴增——熵增；水蒸气在镜上凝成雾，气变液——熵减；小苏打受热放出 $\chem{CO_2}$，气体分子从无到有——熵增；植物把 $\chem{CO_2}$ 和水组装成葡萄糖，自由小分子被锁进有序大分子——熵减。最后一例最易误判：它熵减却不违反第二定律，因为太阳那一端的熵增远远更大——熵铁律只管孤立体系的总账。
\item （“纸可燃、$\Delta G<0$，所以这本书随时会自己烧起来”错在哪？）\textbf{提示：}把“能不能”和“动不动”拆开。\textbf{思路：}$\Delta G<0$ 只判方向：一旦烧起来会自发朝产物走。但启动前分子要先翻过一座能量高坡，常温下能翻过去的寥寥无几，速率慢到等于零。拦住纸张的那道坎叫\textbf{活化能}，正是第 29 章的主角。方向是热力学问题，快慢是动力学问题，这张纸上两个答案分别是“能”和“几乎不”。
\item （石灰窑为什么烧那么高温，而不是室温下多等？）\textbf{提示：}先判断这个反应的 $\Delta H$ 和 $\Delta S$ 符号。\textbf{思路：}$\chem{CaCO_3} \rightarrow \chem{CaO} + \chem{CO_2}$：吸热（不利）但产生气体（$\Delta S>0$，有利）。只有温度足够高，$T\Delta S$ 才能压过 $\Delta H$——实际要烧到约 $900\,^{\circ}\mathrm{C}$ 以上。“室温多等”毫无用处：低温下 $\Delta G>0$，\textbf{方向本身就不允许}，等一亿年也不会发生。等待只能解决动力学问题（慢），解决不了热力学问题（不行）。
\end{enumerate}

\section{样章四：第 49 章《酶：让正确反应在正确时间发生》}

本章的核心：酶是蛋白质做的催化剂，削低活化能、不改 $\Delta G$、不改平衡，靠口袋形状实现专一性。

\begin{enumerate}[itemsep=3pt]
\item （画有酶和无酶两幅能量图，指出哪根箭头变了。）\textbf{提示：}起点和终点由谁决定？\textbf{思路：}两幅图的起点和终点完全一样，$\Delta G$ 箭头长度不变——它只由始末状态决定，与走哪条路无关。变的是活化能箭头：有酶的山顶明显更矮。结论顺势而出：平衡位置由 $\Delta G$ 决定，酶动不了 $\Delta G$，所以酶\textbf{不可能改变平衡}，它只让你\textbf{更快到达}同一个平衡。
\item （估算猪肝里的过氧化氢酶够不够。）\textbf{提示：}两步乘法。\textbf{思路：}$\ud{500}{g}$ 肝约 $10^{11}$ 个细胞，每个细胞约 $10^{7}$ 个酶分子，总数约 $10^{18}$ 个；处理那瓶双氧水需要约 $10^{16}$ 个。富余约 $100$ 倍——绰绰有余。这就是“一滴肝汁就让双氧水疯狂冒泡”的数量级解释：生命在关键酶上从不吝惜库存。
\item （面团醒发与蒸制，用温度—速率曲线解释。）\textbf{提示：}曲线先升后降，拐点是两种效应的交接。\textbf{思路：}升温前期，热运动加快，翻越活化能坡的次数变多，速率上升；超过最适温度后，蛋白质变性、口袋塌陷，速率断崖式下跌。醒发的温暖环境（约 $30\,^{\circ}\mathrm{C}$ 到 $40\,^{\circ}\mathrm{C}$）落在上升段，酵母的酶全力产气；蒸制时 $100\,^{\circ}\mathrm{C}$ 落在深降段，酶失活、酵母死亡——馒头的蓬松全靠蒸之前攒下的气泡，蒸只是定型。
\item （画竞争性抑制的粒子图，解释为什么能被加底物逆转。）\textbf{提示：}竞争的是同一个口袋。\textbf{思路：}只有底物时，所有缺口都被三角占据；加入抑制剂后，形状相似的方块抢走一部分口袋，那部分酶暂时停工；底物大大增加后，三角靠数量优势把方块挤出去，活性恢复。竞争性抑制可逆转，因为抢的是\textbf{同一个}结合位点，谁多谁赢；非竞争性抑制剂结合在口袋\textbf{之外}，把口袋形状拧坏——底物再多，也修不好变了形的口袋。
\item （胃蛋白酶和胰蛋白酶互换环境会怎样？）\textbf{提示：}口袋的形状靠什么维持？\textbf{思路：}口袋靠侧链间的静电吸引、氢键等精细作用维持，而侧链带不带电由 pH 决定。pH 从约 $2$ 换到约 $8$，一批侧链的带电状态翻转，原来的吸引变排斥，口袋变形，酶失活——胃蛋白酶进了小肠就罢工，胰蛋白酶掉进胃里还会被当成蛋白质消化掉。对“口服酶类保健品”的含义：酶是蛋白质，要活着穿过胃酸几乎不可能；判断这类产品，先用附录 E 的清单问证据。
\item （“超级酶能把反应物 $100\%$ 转化成产物”，逐条反驳。）\textbf{提示：}回忆酶的三条“不为”。\textbf{思路：}其一，酶不改变 $\Delta G$——平衡转化率由热力学决定；其二，酶不改变平衡常数和平衡位置——$100\%$ 转化意味着把平衡整个推到产物端，超出酶的职权；其三，酶对正逆反应一视同仁地加速，只缩短到达平衡的时间。真想提高平衡时的产物比例，正确思路是动 $\Delta G$：把产物不断移走（第 31 章）、改变温度，或与一个放能更多的反应耦联（第 28 章）——细胞正是靠耦联办到这件事的。
\end{enumerate}

\section{样章五：第 69 章《翻译：四种字母怎样变成二十种氨基酸》}

本章的核心：翻译是两种语言之间的换码，tRNA 是活字典，每一步准确性都有守门人，每一步都要付能量账。

\begin{enumerate}[itemsep=3pt]
\item （画从 DNA 到成品蛋白质的信息流图。）\textbf{提示：}先写主线，再把“谁执行、说什么语言、能量花在哪”挂到每个箭头上。\textbf{思路：}主线：DNA $\rightarrow$ mRNA $\rightarrow$ 肽链 $\rightarrow$ 折叠修饰后的蛋白质。转录箭头旁：仍是核苷酸语言，执行者 RNA 聚合酶；翻译箭头旁：换成氨基酸语言，执行者核糖体，tRNA 认字扛货，\textbf{氨酰-tRNA 合成酶负责装货，这一步耗 ATP，是准确性的第一道关口}；肽链形成后还有三站：折叠（分子伴侣）、修饰（剪切、加糖链、磷酸化）、运输（信号肽当地址）。对照信息流图清单：换码标在翻译了吗？能量和执行者标了吗？
\item （翻译 AUGUAUGGCUAA，再把阅读框挪一位。）\textbf{提示：}三个字母一组，从 $5'$ 端起读，遇到终止子收工。\textbf{思路：}分组为 AUG--UAU--GGC--UAA：甲硫氨酸（AUG 兼作起始信号）、酪氨酸、甘氨酸，UAA 是终止密码子，到此为止（具体对应以本章密码子片段表为准）。阅读框向后挪一位，变成 UGU--AUG--GCU，读出完全不同的序列，终止子可能消失——同一串字母，换个起点就是另一篇文章。这就是移码突变破坏力巨大的原因：它改的不是一个词，是出错位置往后的整段话。
\item （算一条 500 个氨基酸的多肽要多久、多少 ATP 当量。）\textbf{提示：}两步乘法。\textbf{思路：}每秒约加 20 个氨基酸，500 个要约 $500\div 20=25$ 秒；每个消耗 4 个 ATP 当量，共约 $2000$ 个。再想一层：一个细胞有几千种蛋白质、每种成千上万份——蛋白质合成是细胞最大的能耗项目之一，这就是第 52、53 章那套供能系统片刻不能停的原因。
\item （外星生命用 5 种字母、2 个一组编码 20 种氨基酸，够用吗？）\textbf{提示：}算组合总数，再比冗余。\textbf{思路：}$5^2=25$ 种组合，大于 20，理论上够用。但我们的方案 $4^3=64$，密码子数是氨基酸数的 3 倍多，多出来的冗余叫简并——好几种密码子对应同一种氨基酸，点突变常常“白变”，等于自带纠错缓冲。外星方案只有 5 个空位，几乎每个突变都会改写氨基酸，抗错能力差得多；好处是更省材料。我们的方案是用信息长度换稳健性——这笔交易，生命在三十八亿年前就做了。
\item （红霉素凭什么只卡细菌不卡你？长期吃广谱抗生素为什么肠胃不适？）\textbf{提示：}前半看结构，后半看第 62 章。\textbf{思路：}细菌核糖体和人体细胞质的核糖体在进化上分家已久，形状细节不同——红霉素的“卡口”只匹配细菌版本。药物的选择性毒性，本质都是结构差异。但肠道里住着上万亿共生细菌，它们也是细菌——广谱抗生素不分敌我地轰炸，菌群生态系统被掀翻，消化和屏障功能都受牵连。这是“我们其实是一个生态系统”最直白的一课。
\item （设计实验验证“AAA 编码赖氨酸”。）\textbf{提示：}给体系一封只有一个字母的信，看它译出什么。\textbf{思路：}人工合成只含腺嘌呤的 mRNA（poly-A），放进无细胞翻译体系——含核糖体、tRNA、氨酰-tRNA 合成酶、能量和全部二十种氨基酸，其中赖氨酸做上放射性标记。若产物是清一色的聚赖氨酸，就证明 AAA 对应赖氨酸。对照：换 poly-C 看产物是否换成另一种均聚物；不加 mRNA 看是否没有产物。这正是 1961 年尼伦伯格破译第一格密码表的真实思路——好实验的标准就三条：变量单一、对照齐全、结果可被同行重复。
\end{enumerate}

\section{样章六：第 73 章《从基因到性状，中间隔着整个生命》}

本章的核心：基因到性状是一条多环节的链，每环都可能被环境和其他基因改写；遗传率是群体统计量，不是个人配方。

\begin{enumerate}[itemsep=3pt]
\item （画从 $\mathrm{Hb^{S}}$ 等位基因到疼痛的信息流图。）\textbf{提示：}每环问“谁在驱动”，再问“哪环最容易被撬动”。\textbf{思路：}链条：$\mathrm{Hb^{S}}$ 等位基因 $\rightarrow$ mRNA $\rightarrow$ $\beta$ 珠蛋白第 6 位谷氨酸换成缬氨酸 $\rightarrow$ 脱氧时血红蛋白聚成纤维 $\rightarrow$ 红细胞镰变 $\rightarrow$ 堵塞毛细血管、破裂 $\rightarrow$ 缺氧与疼痛。各环驱动力依次是：碱基互补（信息层）、侧链的疏水相互作用（化学层）、红细胞力学性质（物理层）、血流与供氧（生理层）。最易被环境改变的环在\textbf{中下游}：脱水、感染、高原缺氧会加剧聚合，补水、供氧、医疗可以打断链条——基因在最上游，影响力却要在每一环被重新谈判。
\item （“身高遗传率 $80\%$，所以我的一米七五里一米四归基因”错在哪？）\textbf{提示：}两处错误：一处用错对象，一处用错算术。\textbf{思路：}错误一，遗传率是\textbf{群体}概念：它说“这群人身高差异的约 $80\%$ 与基因差异相关”，对任何单独的个人说“你的身高几分之几归基因”是无意义的。错误二，身高不是“基因块加环境块”的简单加法，基因与环境在发育中持续互作，切不成两块干净分账。两块田的解释：同一批种子分别种在肥沃田和贫瘠田，田\textbf{内部}的高矮差异几乎全来自基因（遗传率接近 $100\%$），而两田之间的巨大平均差几乎全来自环境——遗传率再高，也挡不住环境改天换地。
\item （同样纯合 $\mathrm{Hb^{S}Hb^{S}}$，为什么症状轻重悬殊？）\textbf{提示：}外显率管“演不演”，表现度管“演多狠”。\textbf{思路：}外显率接近 $100\%$ 只保证“会发病”，不保证严重程度。链条上至少两个环节能拉开差距：一是\textbf{其他基因}——胎儿血红蛋白（HbF）水平高的患者，异常血红蛋白被稀释，聚合减少，症状明显轻，而 HbF 高低由别的基因座决定；二是\textbf{环境与医疗}——饮水、感染控制、规范治疗。基因型相同，但基因型只是剧本，同台演出的还有整个基因组、整个身体和整个生活环境。
\item （父母都会卷舌、孩子不会，有哪些解释？）\textbf{提示：}先怀疑“一对基因决定”这个前提本身。\textbf{思路：}至少四种解释：其一，卷舌不是单基因性状，多个基因共同参与，父母的基因恰好在孩子身上凑出“不会”的搭配；其二，外显率不全；其三，环境成分——卷舌涉及舌肌的发育和练习，有人带了对的基因型却从未练出来；其四，统计误差或新发变异。教科书那句“一对等位基因决定”是过度简化，你的调查反例不是 bug，是证据。
\item （人群 A 上训练的多基因风险评分，用于人群 B 为什么失准？）\textbf{提示：}评分学到的是“哪群人的什么规律”？\textbf{思路：}风险评分本质是一台统计机器，它记住的是人群 A 的三样东西：等位基因频率、连锁结构（哪些位点“搭车”）、环境分布。人群 B 这三样全不同：频率变了，权重就错；连锁变了，被当作“标记”的位点可能不再和真凶位点搭同一班车；环境变了，同样基因型的效应大小也变。评分从来不是“基因说明书”，而是“某个人群的统计快照”——换人群，等于拿别人的地图走自己的路。
\item （疟疾被消灭后，$\mathrm{Hb^{S}}$ 频率往哪变？）\textbf{提示：}把选择的天平重新称一遍：杂合优势还在吗？\textbf{思路：}疟疾流行区里，杂合子抗疟疾又不发病，是选择的最爱，$\mathrm{Hb^{S}}$ 因此被维持；疟疾一旦消失，杂合优势归零，而纯合致病的代价还在——天平反转，$\mathrm{Hb^{S}}$ 会被选择缓慢往下压。注意“缓慢”：每代只有纯合子被淘汰，藏在杂合子里的 $\mathrm{Hb^{S}}$ 不暴露，几代之内只见轻微下降。方向是降，速度以百年计。完整的定量工具在第 76 章——这里先练“环境一变，选择的账就要重算”这个直觉。
\end{enumerate}

\section{十一篇的进一步阅读}

下面按篇推荐，每篇三到四条。标准只有三个：内容靠谱、适合青少年、容易找到。书目以中文版或中文字幕为准，具体版本请你自行在图书馆或书店核实；纪录片建议在家长或老师帮助下选择正版渠道。每条后面注明它更适合谁——\textbf{（故事线）}适合想接着追主线的读者，\textbf{（桥梁）}适合想把定量规律抠细的读者，\textbf{（挑战）}适合想啃硬骨头的读者。

\subsection*{第一篇　物质侦探的工具箱}

\begin{itemize}[itemsep=2pt]
\item 《视觉之旅：神奇的化学元素》（西奥多·格雷）：每种元素一页实物照片和故事，把元素从符号变回看得见的东西。（故事线）
\item PhET 互动模拟（美国科罗拉多大学开发，免费，有中文界面）中关于物态、密度和测量的模块：动手拖一拖，比读三遍定义管用。（桥梁）
\item 《元素的故事》（依·尼查叶夫）：老一辈科普经典，讲化学家怎样在瓶瓶罐罐里抓出证据。（故事线）
\item 中国科技馆及各省市科技馆的化学、材料展区：很多展品就是第一篇的实验放大版。（故事线）
\end{itemize}

\subsection*{第二篇　原子与周期表：元素为什么有性格}

\begin{itemize}[itemsep=2pt]
\item 《上帝掷骰子吗：量子物理史话》（曹天元）：量子力学发现史，写得像侦探小说，可跳过公式读故事。（挑战）
\item BBC 纪录片《化学史》（Chemistry: A Volatile History）：从炼金术到现代化学的真实曲折。（故事线）
\item PhET 的“构建一个原子”模拟：亲手添质子、排电子，周期表的形状会自己长出来。（桥梁）
\item 各地科技馆的化学厅或元素周期表主题展：找一找真实的单质样品。（故事线）
\end{itemize}

\subsection*{第三篇　化学键与分子结构：原子为什么结伴}

\begin{itemize}[itemsep=2pt]
\item PhET 的“构建一个分子”和“分子形状”模拟：拖原子、拼分子、转三维模型，形状感是这样练出来的。（桥梁）
\item 《视觉之旅：化学世界的分子奥秘》（西奥多·格雷）：元素那本的下集，讲分子和材料。（故事线）
\item MolView 等免费在线分子模型网站：输入名称就能看到真实分子的三维结构——理解“形状决定功能”的最好玩具。（挑战）
\item 地质博物馆的矿物晶体展厅：天然晶体的规则外形，就是粒子排列写给肉眼看的信。（故事线）
\end{itemize}

\subsection*{第四篇　化学反应：物质怎样重新组合}

\begin{itemize}[itemsep=2pt]
\item PhET 的酸碱、反应速率与可逆反应模拟：把温度、浓度、催化剂旋钮各拧一遍。（桥梁）
\item 正规出版的家庭安全实验类图书（选明确标注“儿童安全实验”的品种，遵守书中安全提示）：厨房酸碱指示剂、晶体培养都值得亲手做。（故事线）
\item 碱金属与水反应的教学演示视频（只许看视频或老师演示，严禁在家模仿）：把“越来越暴躁”看一遍，胜过背十遍。（故事线）
\item 化工博物馆或大型科技馆的工业化学展区：合成氨、炼钢的巨型装置，就是化学方程式放大一亿倍的样子。（故事线）
\end{itemize}

\subsection*{第五篇　有机化学：碳搭出的分子世界}

\begin{itemize}[itemsep=2pt]
\item MolView 里搜一搜阿司匹林、咖啡因、青蒿素：看清官能团长在哪儿，比背结构式有用。（挑战）
\item 科技馆或博物馆的石油、材料、纺织主题展：塑料和纤维展柜前想一想第 43 章的聚合反应。（故事线）
\item 关于香水、酿酒、酿醋工艺的纪录片或专题博物馆：第五篇的分子都藏在香气和酸味里。（故事线）
\end{itemize}

\subsection*{第六篇　生物化学：分子怎样组成生命}

\begin{itemize}[itemsep=2pt]
\item BBC 纪录片《人体奥秘》（Inside the Human Body）：把细胞、代谢、心跳拍给你看。（故事线）
\item 《细胞生命的礼赞》（刘易斯·托马斯）：一位医生关于细胞的散文集，个别篇目偏深，挑能读进去的读。（挑战）
\item Foldit 蛋白质折叠解谜游戏（免费电脑游戏）：亲手把肽链折成低能量形状，第 48、49 章立刻变成手感。（挑战）
\item 动画《工作细胞》及其漫画（中文版已引进）：拟人化讲分子与免疫——记住，拟人是比喻，分子没有性格，只有能量与概率。（故事线）
\end{itemize}

\subsection*{第七篇　细胞：会维护自己的化学城市}

\begin{itemize}[itemsep=2pt]
\item 《工作细胞》系列：免疫细胞各兵种的分工，和第七篇对应得很好。（故事线）
\item 《病毒星球》（卡尔·齐默）：薄薄一本，讲病毒这种“算不上完整生命的化学包裹”怎样塑造生物圈。（故事线）
\item 学校实验室或科普机构的显微镜观察活动（洋葱表皮、草履虫）：第一次亲眼看到细胞的感觉，照片给不了。（桥梁）
\end{itemize}

\subsection*{第八篇　从性状到 DNA}

\begin{itemize}[itemsep=2pt]
\item 《双螺旋》（詹姆斯·沃森）：DNA 结构发现的第一人称记录；作者有争议，书中观点不等于全部史实，可对照第 65 章读。（故事线）
\item 孟德尔与豌豆的传记类科普（图书馆人物传记区有多种青少年版本）：一个修道士怎样用八年豌豆种植发明遗传学。（故事线）
\item PhET 的基因表达基础模拟：亲手把一段 DNA 转录、翻译成蛋白质，为第九篇热身。（桥梁）
\end{itemize}

\subsection*{第九篇　基因如何工作}

\begin{itemize}[itemsep=2pt]
\item PhET 基因表达模拟的进阶玩法：拨动调控开关，看蛋白质产量怎么变——第 70 章的核心直觉。（桥梁）
\item 纪录片《人类本性》（Human Nature，2019，CRISPR 主题）：科学家、患者和伦理学家同场讲述基因编辑。（挑战）
\item 《基因传》（悉达多·穆克吉）：从孟德尔到基因组计划的宏大叙事，篇幅不小，可只读感兴趣的章节。（挑战）
\end{itemize}

\subsection*{第十篇　进化：基因在时间中的故事}

\begin{itemize}[itemsep=2pt]
\item 大卫·爱登堡解说的自然纪录片系列（如《生命的进化》《蓝色星球》）：演化的证据就活在画面里。（故事线）
\item 《地球上最伟大的表演：进化的证据》（理查德·道金斯）：把证据链一条条摆出来，文风犀利，可以跳读。（挑战）
\item 自然博物馆的化石展厅（如国家自然博物馆、上海自然博物馆）：亲眼看到过渡形态的骨架，胜过千言万语。（故事线）
\item PhET 的自然选择模拟：调整环境和捕食压力，看种群频率一代代改变——选择不是故事，是可计算的过程。（桥梁）
\end{itemize}

\subsection*{第十一篇　读写生命与共同未来}

\begin{itemize}[itemsep=2pt]
\item 纪录片《人类本性》（Human Nature）：第九篇没看的话，这一篇务必补上——CRISPR 的希望与边界都在里面。（故事线）
\item 《基因传》（悉达多·穆克吉）后半部分：基因治疗与基因组时代的真实故事。（挑战）
\item 关于转基因、疫苗、基因检测的公开争论：不要急着站队，先用附录 E 的六条清单把双方证据各过一遍筛子。（桥梁）
\item 科技馆、自然博物馆的生态与环境展区：把第 86 章的“生态系统也是化学与信息网络”对号入座。（故事线）
\end{itemize}

\section*{最后一句}

如果这套清单你用熟了，这个附录就完成了使命——因为到那时，答案在谁手里已经不重要了。合上书，找一张纸，随便挑一样你身边的东西（一粒盐、一片叶子、一粒药片），画一画它的粒子图、能量图、物质流图或信息流图，然后自己给自己判卷。能这样做题的人，已经不需要烧掉课本——他早就把课本烧成了自己的问题。
